\documentclass[11pt,a4paper]{scrartcl}

\usepackage{../manostilius_lt}

\title{Kvadratinio daugianario ekstremumai}
\author{Andrius Berniukevičius}

\begin{document}

\maketitle

\section{Pilnojo kvadrato reikšmė}

Norėdami rasti kvadratinio trinario didžiausią arba mažiausią reikšmę, jį pertvarkome išskirdami pilnąjį kvadratą. Pagrindinė idėja yra tokia: bet kurio reiškinio kvadratas negali būti neigiamas:
\[
(x-m)^2\geq0.
\]
Mažiausia kvadrato reikšmė yra $0$, ir ji pasiekiama tada, kai $x=m$.

\begin{definition}
\vocab{Kvadratinio trinario mažiausia reikšmė} -- tai mažiausia reikšmė, kurią gali įgyti reiškinys, kai $x$ įgyja visas galimas reikšmes.
\end{definition}

\begin{definition}
\vocab{Kvadratinio trinario didžiausia reikšmė} -- tai didžiausia reikšmė, kurią gali įgyti reiškinys.
\end{definition}

\begin{definition}
\vocab{Kvadratinio daugianario ekstremumai} -- tai jo mažiausia arba didžiausia reikšmė.
\end{definition}

Kvadratinis daugianaris gali turėti tik vieną iš šių ekstremumų: mažiausią reikšmę arba didžiausią reikšmę. Išskyrę pilnąjį kvadratą galime nustatyti, kuri iš jų yra pasiekiama ir su kuria $x$ reikšme.

\section{Paprasčiausias atvejis: \texorpdfstring{$ax^2$}{ax2}}

Pirmiausia panagrinėkime reiškinį $ax^2$. Kadangi $x^2\geq0$, jo mažiausia reikšmė yra $0$, pasiekiama tada, kai $x=0$. Toliau svarbus tik skaičiaus $a$ ženklas.

Jei $a>0$, padauginę neneigiamą skaičių $x^2$ iš teigiamo skaičiaus gauname neneigiamą reiškinį:
\[
ax^2\geq0.
\]
Todėl reiškinio $ax^2$ mažiausia reikšmė yra $0$, ji pasiekiama, kai $x=0$.

Jei $a<0$, padauginę neneigiamą skaičių $x^2$ iš neigiamo skaičiaus gauname neigiamą arba nulinį reiškinį:
\[
ax^2\leq0.
\]
Todėl reiškinio $ax^2$ didžiausia reikšmė yra $0$, ji pasiekiama, kai $x=0$.

\begin{example}
Raskite reiškinio $3x^2$ ekstremumą.
\textbf{Sprendimas}\\
Kadangi $3>0$ ir $x^2\geq0$, turime $3x^2\geq0$. Mažiausia reikšmė $0$ gaunama, kai $x=0$.
\textbf{Atsakymas:} mažiausia reikšmė yra $0$, ji pasiekiama, kai $x=0$.
\end{example}

\begin{example}
Raskite reiškinio $-5x^2$ ekstremumą.
\textbf{Sprendimas}\\
Kadangi $-5<0$ ir $x^2\geq0$, turime $-5x^2\leq0$. Didžiausia reikšmė $0$ gaunama, kai $x=0$.
\textbf{Atsakymas:} didžiausia reikšmė yra $0$, ji pasiekiama, kai $x=0$.
\end{example}

\section{Mažiausios reikšmės radimas}

Jeigu pilnąjį kvadratą išskyrę gauname formą
\[
a(x-m)^2+n, \qquad a>0,
\]
tai reiškinio mažiausia reikšmė yra $n$. Ji pasiekiama, kai $x=m$.

\begin{property}[Mažiausia reikšmė]
Jei reiškinį galime užrašyti forma $a(x-m)^2+n$, kai $a>0$, tai jo mažiausia reikšmė yra $n$. Ji pasiekiama, kai $x=m$.
\end{property}

Taip yra todėl, kad $a(x-m)^2\geq0$. Vadinasi, prie skaičiaus $n$ visada pridedame neneigiamą skaičių; mažiausiai gauname, kai jis lygus nuliui.

\begin{example}
Raskite reiškinio mažiausią reikšmę ir $x$ reikšmę, su kuria ji pasiekiama:
\[
x^2+6x+2.
\]
\textbf{Sprendimas}\\
Kadangi $6x=2\cdot x\cdot3$, pridedame ir atimame $3^2=9$:
\[
\begin{aligned}
x^2+6x+2
&=x^2+6x+9-9+2\\
&=(x+3)^2-7.
\end{aligned}
\]
Kadangi $(x+3)^2\geq0$, mažiausia jo reikšmė yra $0$. Ji gaunama, kai $x+3=0$, t. y. kai $x=-3$. Todėl $(x+3)^2-7\geq-7$.\\ \\
\textbf{Atsakymas:} mažiausia reikšmė yra $-7$, ji pasiekiama, kai $x=-3$.
\end{example}

\begin{example}
Raskite reiškinio mažiausią reikšmę ir $x$ reikšmę, su kuria ji pasiekiama:
\[
2x^2-8x+3.
\]
\textbf{Sprendimas}\\
Iškeliame $2$ prie pirmųjų dviejų narių ir išskiriame pilnąjį kvadratą:
\[
\begin{aligned}
2x^2-8x+3
&=2(x^2-4x)+3\\
&=2(x^2-4x+4-4)+3\\
&=2(x-2)^2-5.
\end{aligned}
\]
Kadangi $2(x-2)^2\geq0$, mažiausia reikšmė yra $-5$. Ji pasiekiama, kai $x-2=0$, t. y. kai $x=2$.\\ \\
\textbf{Atsakymas:} mažiausia reikšmė yra $-5$, ji pasiekiama, kai $x=2$.
\end{example}

\section{Didžiausios reikšmės radimas}

Jeigu pilnąjį kvadratą išskyrę gauname formą
\[
-a(x-m)^2+n, \qquad a>0,
\]
tai reiškinio didžiausia reikšmė yra $n$. Ji pasiekiama, kai $x=m$.

\begin{property}[Didžiausia reikšmė]
Jei reiškinį galime užrašyti forma $-a(x-m)^2+n$, kai $a>0$, tai jo didžiausia reikšmė yra $n$. Ji pasiekiama, kai $x=m$.
\end{property}

Taip yra todėl, kad $-a(x-m)^2\leq0$. Vadinasi, prie skaičiaus $n$ pridedame neigiamą arba nulinį skaičių; daugiausiai gauname, kai jis lygus nuliui.

\begin{example}
Raskite reiškinio didžiausią reikšmę ir $x$ reikšmę, su kuria ji pasiekiama:
\[
-x^2+4x+1.
\]
\textbf{Sprendimas}\\
Iškeliame minusą prie pirmųjų dviejų narių ir išskiriame pilnąjį kvadratą:
\[
\begin{aligned}
-x^2+4x+1
&=-(x^2-4x)+1\\
&=-(x^2-4x+4-4)+1\\
&=-(x-2)^2+5.
\end{aligned}
\]
Kadangi $-(x-2)^2\leq0$, didžiausia reikšmė yra $5$. Ji pasiekiama, kai $x-2=0$, t. y. kai $x=2$.\\ \\
\textbf{Atsakymas:} didžiausia reikšmė yra $5$, ji pasiekiama, kai $x=2$.
\end{example}

\begin{example}
Raskite reiškinio didžiausią reikšmę ir $x$ reikšmę, su kuria ji pasiekiama:
\[
-3x^2-12x+7.
\]
\textbf{Sprendimas}\\
Iškeliame $-3$ prie pirmųjų dviejų narių:
\[
\begin{aligned}
-3x^2-12x+7
&=-3(x^2+4x)+7\\
&=-3(x^2+4x+4-4)+7\\
&=-3(x+2)^2+19.
\end{aligned}
\]
Kadangi $-3(x+2)^2\leq0$, didžiausia reikšmė yra $19$. Ji pasiekiama, kai $x+2=0$, t. y. kai $x=-2$.\\ \\
\textbf{Atsakymas:} didžiausia reikšmė yra $19$, ji pasiekiama, kai $x=-2$.
\end{example}

\begin{remark}
Išskyrus pilnąjį kvadratą, sprendimo kryptį nusako koeficiento prie kvadrato ženklas:
\[
a(x-m)^2+n,\quad a>0 \quad \Rightarrow \quad \text{mažiausia reikšmė yra } n;
\]
\[
-a(x-m)^2+n,\quad a>0 \quad \Rightarrow \quad \text{didžiausia reikšmė yra } n.
\]
Abiem atvejais ekstremumo reikšmė pasiekiama tada, kai kvadratas lygus nuliui: $x=m$.
\end{remark}

\section{Uždaviniai}

\begin{problem}
Išskirdami pilnąjį kvadratą raskite reiškinio mažiausią arba didžiausią reikšmę ir nurodykite, su kuria $x$ reikšme ji pasiekiama.
\begin{tasks}[label={(\arabic*)}, label-offset=0.9em](2)
	\task $x^2-10x+4$
	\task $x^2+2x-8$
	\task $-x^2-6x+2$
	\task $-2x^2+8x+1$
	\task $3x^2-6x+7$
	\task $-4x^2-16x+3$
\end{tasks}
\end{problem}

\begin{problem}
Išskirdami pilnąjį kvadratą raskite reiškinio mažiausią arba didžiausią reikšmę ir nurodykite, su kuria $x$ reikšme ji pasiekiama.
\begin{tasks}[label={(\arabic*)}, label-offset=0.9em](2)
	\task $2x^2+12x-1$
	\task $5x^2-20x+6$
	\task $-2x^2+4x+9$
	\task $-3x^2+18x-4$
	\task $4x^2+4x+10$
	\task $-x^2+10x-12$
\end{tasks}
\end{problem}

\newpage
\answerssection

\begin{problem} \phantom{text}
\begin{tasks}[label={(\arabic*)}, label-offset=0.9em](2)
	\task mažiausia reikšmė $-21$, kai $x=5$;
	\task mažiausia reikšmė $-9$, kai $x=-1$;
	\task didžiausia reikšmė $11$, kai $x=-3$;
	\task didžiausia reikšmė $9$, kai $x=2$;
	\task mažiausia reikšmė $4$, kai $x=1$;
	\task didžiausia reikšmė $19$, kai $x=-2$.
\end{tasks}
\end{problem}

\begin{problem} \phantom{text}
\begin{tasks}[label={(\arabic*)}, label-offset=0.9em](2)
	\task mažiausia reikšmė $-19$, kai $x=-3$;
	\task mažiausia reikšmė $-14$, kai $x=2$;
	\task didžiausia reikšmė $11$, kai $x=1$;
	\task didžiausia reikšmė $23$, kai $x=3$;
	\task mažiausia reikšmė $9$, kai $x=-\frac{1}{2}$;
	\task didžiausia reikšmė $13$, kai $x=5$.
\end{tasks}
\end{problem}

\end{document}
